%
% Configuratie
%

% Preambule met standaardinstellingen
\documentclass[a4paper,oneside,11pt,final]{memoir}

% Gesplitst document
\usepackage{subfiles}

% Noot: zorg ervoor dat Nederlandse woord-splitsing geactiveerd is.
\usepackage[dutch]{babel}

% UTF8 gebruiken voor gebruik van alle symbolen
\usepackage[utf8]{inputenc}
\usepackage{eurosym}
\DeclareUnicodeCharacter{03A9}{\Omega}

% Noot: je kan het graphicxpakket een optie dvips of pdftex doorgeven
% in dat geval moet je ze ook aan iiiscriptie doorgeven, dus bijvoorbeeld
% \usepackage[dvips]{graphicx}
% \usepackage[dvips]{iiiscriptie}
\usepackage{graphicx}
\usepackage{iiiscriptie}

% Tabellen eleganter maken
\usepackage{booktabs}

% Navigeerbaarheid van hyperlinks in PDF
\usepackage{hyperref}

% Afkortingen intelligent gebruiken
\usepackage[printonlyused,withpage]{acronym}

% Math support
\usepackage{amssymb,amsmath}

% Extra functies
% Verkleinde margin entry
\setlength{\marginparwidth}{1.2in}
\let\oldmarginpar\marginpar
\renewcommand\marginpar[1] {\-\oldmarginpar[\raggedleft\footnotesize #1]%
{\raggedright\footnotesize #1}}

% Een TODO-entry
\newcommand{\todo}[1] {
	\addcontentsline{tdo}{todo}{\protect{#1}}
	\marginpar{#1}
}

% Hyperlink maken en URL in footnote tonen
\usepackage{hyperref}
\newcommand{\makeurl}[2]{\href{#1}{#2} \footnote{#1}}

% Compacte enumeraties
\newenvironment{enumerate_compact}{
\begin{enumerate}
  \setlength{\itemsep}{1pt}
  \setlength{\parskip}{0pt}
  \setlength{\parsep}{0pt}
}{\end{enumerate}}
\newenvironment{itemize_compact}{
\begin{itemize}
  \setlength{\itemsep}{1pt}
  \setlength{\parskip}{0pt}
  \setlength{\parsep}{0pt}
}{\end{itemize}}




%
% Titelpagina
%

% Invullen velden
\departement{Departement Toegepaste Ingenieurswetenschappen}
\deptadres{Valentin Vaerwyckweg 1, 9000 Gent}
\studiejaar{1e Master Informatica}
\soortrapport{
Scriptie voorgedragen tot het behalen van de diploma van\\
MASTER IN DE TOEGEPASTE INGENIEURSWETENSCHAPPEN: INFORMATICA\\
}
\title{Ontwikkeling van een multimediaframework}
\bedrijfslogo{\includegraphics[height=18mm]{logomira}}
\author{Tim BESARD}
\promotoren{
Intern:& Leen BROUNS \\
Extern:& Philippe MOLLET}


%
% Inhoud
%

\begin{document}

\pagenumbering{roman}
%
% Titelpagina
%

\maketitle


%
% Abstract
%

\chapter*{Abstract}
\addcontentsline{toc}{chapter}{Abstract}

\textit{Hier komt een abstract.}


%
% Voorwoord
%

\chapter*{Voorwoord}
\addcontentsline{toc}{chapter}{Voorwoord}

\textit{Hier komt het voorwoord.}


%
% Inhoudstafel
%

\setlength\cftpartnumwidth{2em}

\newpage
\tableofcontents

\newpage


%
% Lijst met afbeeldingen
%

\listoffigures


%
% Lijst van afkortingen
%

\chapter*{Lijst van afkortingen}
\addcontentsline{toc}{chapter}{Lijst van afkortingen}

\begin{acronym}[WYSIWYG]	% langste afkorting

\acro{upnp}[UPnP]{Universal Plug and Play}
\acro{mdns}[mDNS]{Multicast DNS}
\acro{dns}[DNS]{Domain Name System}
\acro{ssdp}[SSDP]{Simple Service Discovery Protocol}
\acro{dcp}[DCP]{Device Control Protocol}
\acro{ssl}[SSL]{Secure Socket Layer}
\acro{rpc}[RPC]{Remote Procedure Call}
\acro{webdav}[WebDAV]{Web-based Distributed Authoring and Versioning}
\acro{http}[HTTP]{Hypertext Transport Protocol}
\acro{rmi}[RMI]{Remote Method Invocation}
\acro{corba}[CORBA]{Common Object Request Broker Architecture}
\acro{soap}[SOAP]{Simple Object Access Protocol}
\acro{wysiwyg}[WYSIWYG]{What You See Is What You Get}
\acro{html}[HTML]{Hypertext Markup Language}
\acro{rest}[REST]{Representational State Transfer}
\acro{jvm}[JVM]{Java Virtual Machine}
\acro{svn}[SVN]{Subversion}
\acro{ram}[RAM]{Random Access Memory}
\acro{nas}[NAS]{Network Attached Storage}
\acro{gpu}[GPU]{Graphics Processing Unit}
\acro{dsp}[DSP]{Digital Signal Processor}
\acro{usb}[USB]{Universal Serial Bus}
\acro{hid}[HID]{Human Interface Device}
\acro{lufa}[LUFA]{Lightweight USB Framework for AVRs}
\acro{pll}[PLL]{Phase-Locked Loop

\end{acronym}

\newpage

\pagenumbering{arabic}
\subfile{1_introductie.tex}
\subfile{2_ontwerp.tex}
\subfile{3_server.tex}
\subfile{4_kiosk.tex}
\subfile{5_inputmodule.tex}
\subfile{6_invoering.tex}

\chapter{Conclusie}

% todo: mag niet teveel op abstract lijken
% http://www.dissertationswriting.co.uk/dissertation_conclusion.htm 

Het doel van deze thesis was de ontwikkeling van een modern multimediaframework dat zorgt voor de opslag, distributie en weergave van multimediafragmenten op de kiosken van de volkssterrenwacht MIRA. Daarbij is belangrijk dat een beheerder elk van die processen voldoende kan opvolgen en beïnvloeden, zonder al teveel technische kennis nodig te hebben. Ook moet het mogelijk zijn om mooie voorstellingen te realiseren met de middelen die het platform aanbiedt, en tegelijk compatibel te zijn met de voorstellingen die momenteel gebruikt worden.

We hebben ervoor gekozen om de voorstellingen op te maken in \ac{html} en Javascript, waarmeer we zeer rijke gebruikerservaringen kunnen ontwerpen. Ook compatibiliteit met de huidige voorstellingen is ermee mogelijk: na conversie naar het WebM formaat kunnen we de fragmenten gebruiken binnen \ac{html}. Om daarbij dynamisch de taalkeuze van de gebruiker in te verwerken hebben we voorzien in een kleine webapplicatie, opgebouwd met behulp van JQuery.
Vervolgens worden de voorstellingen opgeslaan in een \ac{svn} repository, waardoor efficiënte overdracht bekomen wordt en een beheerder relatief eenvoudig wijzigingen kan doorvoeren aan de inhoud van de voorstellingen.

% server

% client

% inputmodule

% conclusie

% Bibliografie
\bibliographystyle{plainnat}
\bibliography{verslag}


\end{document}
